% Variante de la Actividad 9 usando java.util.concurrent.Semaphore
\documentclass[12pt,oneside]{template/liverpoolthesis}
% !TEX root = ../Practica1_Hernandez_Andrea_Espino_Heriberto.tex
\usepackage[discipline=computer-science]{template/liverpoolthesis}
\usepackage[utf8]{inputenc}
\usepackage[T1]{fontenc}
\usepackage[spanish,es-nodecimaldot]{babel}

% Las referencias internas y las citas usan el mismo azul que los títulos.
\hypersetup{
  linkcolor = LiverpoolBlue,
  citecolor = LiverpoolBlue,
  urlcolor = LiverpoolBlue
}

% La marca de cada capítulo aparece en el encabezado, arriba de la regla.
\addto\captionsspanish{\renewcommand{\chaptername}{Ejercicio}}
\providecommand{\figureautorefname}{Figura}
\addto\captionsspanish{\renewcommand{\figureautorefname}{Figura}}

\graphicspath{{figuras/}}
\setlength{\headheight}{28pt}

% Inserta una captura de código y/o de ejecución.
% El tercer argumento es la etiqueta; el argumento opcional controla la altura.
\newcommand{\incluirevidencia}[4][0.70\textheight]{%
  \begin{figure}[H]
    \centering
    \includegraphics[
      width=0.95\textwidth,
      height=#1,
      keepaspectratio
    ]{#2}
    \caption{#3}
    \label{#4}
  \end{figure}%
}

% ---------- Datos de la portada ----------
\title{Práctica de laboratorio 1: Creación de procesos e hilos en C y Java}
\author{%
  Heriberto Espino Montelongo\\
  ID: 175199\\[1em]
  Andrea Hernandez Huerta\\
  ID: 175523\\[1em]
  \text{Profesor:} Dr. Victor Giovanni Morales Murillo}
\faculty{Universidad de las Américas Puebla (UDLAP)}
\school{Licenciatura en Ciencia de Datos}
\degree{Curso: Programación Concurrente}
\date{Fecha de entrega: \today}
\subject{Creación y ejecución de procesos e hilos en C y Java}


\title{Actividad 9: Semáforo binario y contador}
\author{%
  Andrea Hernandez Huerta\\
  ID: 175523\\[1em]
  Heriberto Espino Montelongo\\
  ID: 175199\\[1em]
  \text{Profesor:} Dr. Victor Giovanni Morales Murillo}
\faculty{Universidad de las Américas Puebla (UDLAP)}
\school{Licenciatura en Ciencia de Datos}
\degree{Curso: Programación Concurrente}
\date{Fecha de entrega: \today}
\subject{Semáforos con \texttt{java.util.concurrent.Semaphore}}

\addto\captionsspanish{\renewcommand{\chaptername}{Parte}}

\begin{document}

\frontmatter
\maketitle
\tableofcontents
\listoffigures
\listoftables

\mainmatter

\chapter{Semáforo binario: cuenta bancaria}

\textbf{Instrucción.} Crear una cuenta bancaria compartida con saldo inicial de
\$1000 y cinco hilos que realicen depósitos y retiros. El acceso al saldo debe
estar protegido por un semáforo binario con un solo permiso.

\section{Entradas y procesamiento}

El programa utiliza un saldo compartido de \$1000 y un objeto de la clase
\texttt{Semaphore} con un permiso:

\begin{center}
\texttt{private static final Semaphore sem = new Semaphore(1);}
\end{center}

El método \texttt{acquire()} solicita el permiso. Si otro hilo ya lo tiene,
el hilo que intenta entrar queda bloqueado hasta que el permiso sea liberado.
Al terminar la operación sobre el saldo se utiliza \texttt{release()} para
devolver el permiso.

Los cinco hilos realizan tres depósitos de \$200, \$300 y \$50, y dos retiros
de \$150 y \$100. Por lo tanto, el saldo final esperado es

\[
1000+200-150+300-100+50=1300.
\]

En esta implementación, \texttt{acquire()} cumple el papel de la operación
\texttt{P()} estudiada en clase y \texttt{release()} cumple el papel de
\texttt{V()}.

\section{Prueba de escritorio}

Para observar la exclusión mutua se consideran dos hilos: A deposita \$200 y B
retira \$150.

\begin{table}[H]
\centering
\caption{Prueba de escritorio del semáforo binario.}
\begin{tabularx}{\textwidth}{c L c L}
\toprule
\textbf{Paso} & \textbf{Acción} & \textbf{Permisos} & \textbf{Resultado}\\
\midrule
0 & Estado inicial & 1 & El semáforo está disponible.\\
1 & A ejecuta \texttt{acquire()} & 0 & A obtiene el único permiso.\\
2 & B ejecuta \texttt{acquire()} & 0 & B queda esperando.\\
3 & A deposita \$200 & 0 & El saldo cambia de \$1000 a \$1200.\\
4 & A ejecuta \texttt{release()} & 1 & El permiso queda disponible para B.\\
5 & B obtiene el permiso & 0 & B entra y retira \$150.\\
6 & B ejecuta \texttt{release()} & 1 & El semáforo vuelve a quedar disponible.\\
\bottomrule
\end{tabularx}
\end{table}

\section{Evidencia y resultados}

La Figura~\ref{fig:binario} muestra en una sola captura el código utilizado y
varias ejecuciones del programa. El orden de los hilos cambia entre
ejecuciones, pero en todos los casos el saldo final mostrado es \$1300. Esto
confirma que cada hilo obtiene el único permiso antes de modificar la variable
compartida.

\incluirevidencia[0.82\textheight]{1.png}
  {Código y resultados del semáforo binario con \texttt{Semaphore(1)}.}
  {fig:binario}

\chapter{Semáforo contador: impresoras}

\textbf{Instrucción.} Simular tres impresoras compartidas por ocho hilos. El
semáforo debe permitir como máximo tres impresiones simultáneas.

\section{Entradas y procesamiento}

El programa representa las tres impresoras mediante un semáforo con tres
permisos:

\begin{center}
\texttt{private static final Semaphore impresoras = new Semaphore(3);}
\end{center}

Cada hilo ejecuta \texttt{impresoras.acquire()} antes de imprimir. Los tres
primeros hilos que obtienen un permiso pueden continuar; si un cuarto hilo
intenta entrar cuando los tres permisos están ocupados, queda bloqueado.

Después de simular la impresión mediante \texttt{Thread.sleep(1000)}, el hilo
ejecuta \texttt{impresoras.release()} y devuelve su permiso. De esta forma otro
hilo que estaba esperando puede comenzar.

\section{Prueba de escritorio}

Para la prueba se consideran dos permisos y cuatro hilos A, B, C y D.

\begin{table}[H]
\centering
\caption{Prueba de escritorio del semáforo contador.}
\begin{tabularx}{\textwidth}{c L c L}
\toprule
\textbf{Paso} & \textbf{Acción} & \textbf{Permisos} & \textbf{Resultado}\\
\midrule
0 & Estado inicial & 2 & Hay dos impresoras disponibles.\\
1 & A ejecuta \texttt{acquire()} & 1 & A comienza a imprimir.\\
2 & B ejecuta \texttt{acquire()} & 0 & B comienza a imprimir.\\
3 & C ejecuta \texttt{acquire()} & 0 & C queda esperando.\\
4 & D ejecuta \texttt{acquire()} & 0 & D queda esperando.\\
5 & A ejecuta \texttt{release()} & 1 & Uno de los hilos bloqueados puede continuar.\\
6 & C obtiene el permiso & 0 & C comienza a imprimir.\\
7 & B ejecuta \texttt{release()} & 1 & D puede obtener el siguiente permiso.\\
8 & C y D terminan & 2 & Los dos permisos vuelven a estar disponibles.\\
\bottomrule
\end{tabularx}
\end{table}

\section{Evidencia y resultados}

La Figura~\ref{fig:contador} muestra el código y varias ejecuciones en la misma
captura. En la terminal se observa que varios hilos comienzan a imprimir antes
de que terminen los anteriores, pero el semáforo solo dispone de tres permisos.
Por ello, como máximo tres hilos pueden encontrarse usando una impresora al
mismo tiempo. El orden cambia entre ejecuciones porque depende de la
planificación de los hilos.

\incluirevidencia[0.82\textheight]{2.png}
  {Código y resultados del semáforo contador con \texttt{Semaphore(3)}.}
  {fig:contador}

\chapter{Comparación y conclusión}

\section{Comparación}

\begin{table}[H]
\centering
\caption{Comparación entre los dos semáforos utilizados.}
\begin{tabularx}{\textwidth}{P{4.2cm} C C}
\toprule
\textbf{Aspecto} & \textbf{Binario} & \textbf{Contador}\\
\midrule
Creación & \texttt{new Semaphore(1)} & \texttt{new Semaphore(3)}\\
Permisos iniciales & 1 & 3\\
Recurso controlado & Saldo bancario & Tres impresoras\\
Máximo de hilos simultáneos & 1 & 3\\
Solicitar permiso & \texttt{acquire()} & \texttt{acquire()}\\
Liberar permiso & \texttt{release()} & \texttt{release()}\\
Bloqueo & Cuando el permiso está ocupado & Cuando los tres permisos están ocupados\\
\bottomrule
\end{tabularx}
\end{table}

\section{Conclusión}

La clase \texttt{Semaphore} de Java permite aplicar directamente el mecanismo
estudiado en clase. Con un permiso se obtiene exclusión mutua sobre la cuenta
bancaria; con tres permisos se permite que hasta tres hilos utilicen recursos
equivalentes de forma concurrente. En ambos ejercicios,
\texttt{acquire()} solicita un permiso y puede bloquear al hilo, mientras que
\texttt{release()} devuelve el permiso para que otro hilo pueda continuar.

Los resultados observados coinciden con el comportamiento esperado. En la
cuenta bancaria el orden de ejecución puede cambiar, pero el saldo final se
mantiene en \$1300. En el ejercicio de impresoras el orden también varía, pero
la cantidad de hilos que puede avanzar simultáneamente queda limitada por los
tres permisos del semáforo.

\end{document}
